\documentclass[12pt,a4paper]{article}
\usepackage[margin=1in]{geometry}
\usepackage{amsmath,amssymb,amsfonts}
\usepackage{graphicx}
\usepackage{hyperref}
\usepackage{listings}
\usepackage{xcolor}
\usepackage{booktabs}
\usepackage{longtable}
\usepackage{array}
\usepackage{multirow}
\usepackage{tikz}
\usetikzlibrary{shapes,arrows,positioning,fit,calc}
\usepackage{float}
\usepackage{enumitem}
\usepackage{fancyhdr}
\usepackage{tcolorbox}

\pagestyle{fancy}
\fancyhf{}
\rhead{CA Lab 3 -- Theory Document (Part 2)}
\lhead{ToyRISC ISA Specification}
\cfoot{\thepage}

\definecolor{codegreen}{rgb}{0,0.6,0}
\definecolor{codegray}{rgb}{0.5,0.5,0.5}
\definecolor{codepurple}{rgb}{0.58,0,0.82}
\definecolor{backcolour}{rgb}{0.95,0.95,0.92}

\lstdefinestyle{javastyle}{
    backgroundcolor=\color{backcolour},
    commentstyle=\color{codegreen},
    keywordstyle=\color{blue}\bfseries,
    numberstyle=\tiny\color{codegray},
    stringstyle=\color{codepurple},
    basicstyle=\ttfamily\footnotesize,
    breakatwhitespace=false,
    breaklines=true,
    captionpos=b,
    keepspaces=true,
    numbers=left,
    numbersep=5pt,
    showspaces=false,
    showstringspaces=false,
    showtabs=false,
    tabsize=2,
    frame=single
}

\lstdefinestyle{asmstyle}{
    backgroundcolor=\color{backcolour},
    commentstyle=\color{codegreen},
    keywordstyle=\color{blue}\bfseries,
    basicstyle=\ttfamily\footnotesize,
    breaklines=true,
    captionpos=b,
    keepspaces=true,
    numbers=left,
    numbersep=5pt,
    frame=single,
    morekeywords={add,addi,sub,subi,mul,muli,div,divi,and,andi,or,ori,xor,xori,slt,slti,sll,slli,srl,srli,sra,srai,load,store,jmp,beq,bne,blt,bgt,end}
}

\title{\Huge\bfseries Comprehensive Theory of Single-Cycle Processor Design and the ToyRISC ISA\\[1em]
\Large Part 2: ToyRISC ISA Deep Dive and Instruction Encoding}
\author{Computer Architecture Laboratory -- Assignment 3}
\date{\today}

\begin{document}

\maketitle
\tableofcontents
\newpage

% ===========================================================================
\section{ToyRISC ISA Overview}
% ===========================================================================

The ToyRISC ISA is a custom RISC instruction set architecture designed for educational purposes. It is inspired by RISC-V (RV32I) and shares many design principles with it. The ISA is clean, minimal, and regular, making it suitable for building a processor simulator from scratch.

\subsection{Key Characteristics}

\begin{itemize}[leftmargin=*]
    \item \textbf{32-bit fixed-width instructions}: Every instruction is exactly 32 bits, simplifying the fetch and decode logic.
    \item \textbf{32 general-purpose registers}: Named \texttt{x0} through \texttt{x31}, each 32 bits wide.
    \item \textbf{Load-store architecture}: Only \texttt{load} and \texttt{store} access memory. All computations occur between registers.
    \item \textbf{30 distinct operations}: Including arithmetic, logical, memory, control flow, and a special \texttt{end} instruction.
    \item \textbf{3 instruction formats}: R3-Type, R2I-Type, and RI-Type.
    \item \textbf{5-bit opcode}: Allows for up to $2^5 = 32$ unique opcodes (30 are used, leaving 2 reserved).
    \item \textbf{Word-addressable memory}: $2^{16} = 65536$ addressable words.
\end{itemize}

% ===========================================================================
\section{Instruction Formats -- Detailed Encoding}
% ===========================================================================

ToyRISC defines three instruction formats. The opcode is always in the most significant 5 bits (bits 31--27). The remaining 27 bits are interpreted differently based on the format.

\subsection{R3-Type (Register-Register-Register)}

Used by arithmetic/logical instructions that operate on \textbf{three registers} (two sources, one destination).

\begin{center}
\begin{tabular}{|c|c|c|c|c|}
\hline
\textbf{Bits 31--27} & \textbf{Bits 26--22} & \textbf{Bits 21--17} & \textbf{Bits 16--12} & \textbf{Bits 11--0} \\
\hline
opcode (5) & rs1 (5) & rs2 (5) & rd (5) & unused (12) \\
\hline
\end{tabular}
\end{center}

\begin{itemize}
    \item \textbf{opcode} (5 bits, [31:27]): Identifies the operation (e.g., 00000 for \texttt{add}).
    \item \textbf{rs1} (5 bits, [26:22]): First source register number.
    \item \textbf{rs2} (5 bits, [21:17]): Second source register number.
    \item \textbf{rd} (5 bits, [16:12]): Destination register number.
    \item \textbf{unused} (12 bits, [11:0]): Not used; set to 0.
\end{itemize}

\textbf{Example}: \texttt{add \%x3, \%x4, \%x5} (opcode=00000, rs1=00011, rs2=00100, rd=00101)

Binary encoding:
\[
\underbrace{00000}_{\text{add}}\;\underbrace{00011}_{\text{x3}}\;\underbrace{00100}_{\text{x4}}\;\underbrace{00101}_{\text{x5}}\;\underbrace{000000000000}_{\text{unused}}
\]

Decimal: Each field contributes to the 32-bit integer value stored in memory.

\subsubsection{Bit Extraction Formulas for R3-Type}

To decode an R3-Type instruction from a 32-bit integer \texttt{instr}:

\begin{align}
\text{opcode} &= (\texttt{instr} \gg 27) \;\&\; \texttt{0x1F} \\
\text{rs1} &= (\texttt{instr} \gg 22) \;\&\; \texttt{0x1F} \\
\text{rs2} &= (\texttt{instr} \gg 17) \;\&\; \texttt{0x1F} \\
\text{rd} &= (\texttt{instr} \gg 12) \;\&\; \texttt{0x1F}
\end{align}

In Java: \texttt{int opcode = (instr >>> 27) \& 0x1F;}

\subsection{R2I-Type (Register-Register-Immediate)}

Used by arithmetic/logical instructions with an immediate value, memory instructions, and conditional branches.

\begin{center}
\begin{tabular}{|c|c|c|c|}
\hline
\textbf{Bits 31--27} & \textbf{Bits 26--22} & \textbf{Bits 21--17} & \textbf{Bits 16--0} \\
\hline
opcode (5) & rs1 (5) & rd (5) & immediate (17) \\
\hline
\end{tabular}
\end{center}

\begin{itemize}
    \item \textbf{opcode} (5 bits, [31:27]): Identifies the operation.
    \item \textbf{rs1} (5 bits, [26:22]): Source register number.
    \item \textbf{rd} (5 bits, [21:17]): Destination register number.
    \item \textbf{immediate} (17 bits, [16:0]): Signed immediate value (two's complement). Range: $-2^{16}$ to $2^{16} - 1$ (i.e., $-65536$ to $65535$).
\end{itemize}

\textbf{Note}: For R2I-Type, the ISA specification states that immediate values are placed in \texttt{sourceOperand2} in the ParsedProgram representation. This is critical for the Operand Fetch implementation.

\subsubsection{Sign Extension of 17-bit Immediate}

The 17-bit immediate must be sign-extended to 32 bits before use:

\begin{lstlisting}[style=javastyle, caption={Sign extension of 17-bit immediate}]
int imm17 = instr & 0x1FFFF;   // extract lower 17 bits
// Sign extend: if bit 16 is set, the value is negative
if ((imm17 & (1 << 16)) != 0) {
    imm17 = imm17 | 0xFFFE0000; // fill upper 15 bits with 1s
}
\end{lstlisting}

Alternatively, using arithmetic shift:
\begin{lstlisting}[style=javastyle]
int imm17 = (instr << 15) >> 15; // shift left, then arith right
\end{lstlisting}

\subsubsection{Bit Extraction Formulas for R2I-Type}

\begin{align}
\text{opcode} &= (\texttt{instr} \gg 27) \;\&\; \texttt{0x1F} \\
\text{rs1} &= (\texttt{instr} \gg 22) \;\&\; \texttt{0x1F} \\
\text{rd} &= (\texttt{instr} \gg 17) \;\&\; \texttt{0x1F} \\
\text{imm} &= \text{signExtend}_{17 \to 32}(\texttt{instr} \;\&\; \texttt{0x1FFFF})
\end{align}

\subsection{RI-Type (Register-Immediate)}

Used by the \texttt{jmp} and \texttt{end} instructions.

\begin{center}
\begin{tabular}{|c|c|c|}
\hline
\textbf{Bits 31--27} & \textbf{Bits 26--22} & \textbf{Bits 21--0} \\
\hline
opcode (5) & rd (5) & immediate (22) \\
\hline
\end{tabular}
\end{center}

\begin{itemize}
    \item \textbf{opcode} (5 bits, [31:27]): Identifies the operation.
    \item \textbf{rd} (5 bits, [26:22]): Destination register (used by \texttt{jmp}; unused by \texttt{end}).
    \item \textbf{immediate} (22 bits, [21:0]): Signed immediate value. Range: $-2^{21}$ to $2^{21} - 1$.
\end{itemize}

\subsubsection{Bit Extraction Formulas for RI-Type}

\begin{align}
\text{opcode} &= (\texttt{instr} \gg 27) \;\&\; \texttt{0x1F} \\
\text{rd} &= (\texttt{instr} \gg 22) \;\&\; \texttt{0x1F} \\
\text{imm} &= \text{signExtend}_{22 \to 32}(\texttt{instr} \;\&\; \texttt{0x3FFFFF})
\end{align}

\subsection{Opcode-to-Format Mapping}

Not all opcodes use the same format. Here is the complete mapping:

\begin{longtable}{p{1.5cm}p{2cm}p{2cm}p{7cm}}
\toprule
\textbf{Opcode} & \textbf{Binary} & \textbf{Format} & \textbf{Operation} \\
\midrule
0 & 00000 & R3-Type & \texttt{add} rd = rs1 + rs2 \\
1 & 00001 & R2I-Type & \texttt{addi} rd = rs1 + imm \\
2 & 00010 & R3-Type & \texttt{sub} rd = rs1 - rs2 \\
3 & 00011 & R2I-Type & \texttt{subi} rd = rs1 - imm \\
4 & 00100 & R3-Type & \texttt{mul} rd = rs1 $\times$ rs2 \\
5 & 00101 & R2I-Type & \texttt{muli} rd = rs1 $\times$ imm \\
6 & 00110 & R3-Type & \texttt{div} rd = rs1 $\div$ rs2 \\
7 & 00111 & R2I-Type & \texttt{divi} rd = rs1 $\div$ imm \\
8 & 01000 & R3-Type & \texttt{and} rd = rs1 \& rs2 \\
9 & 01001 & R2I-Type & \texttt{andi} rd = rs1 \& imm \\
10 & 01010 & R3-Type & \texttt{or} rd = rs1 $|$ rs2 \\
11 & 01011 & R2I-Type & \texttt{ori} rd = rs1 $|$ imm \\
12 & 01100 & R3-Type & \texttt{xor} rd = rs1 $\oplus$ rs2 \\
13 & 01101 & R2I-Type & \texttt{xori} rd = rs1 $\oplus$ imm \\
14 & 01110 & R3-Type & \texttt{slt} rd = (rs1 $<$ rs2) ? 1 : 0 \\
15 & 01111 & R2I-Type & \texttt{slti} rd = (rs1 $<$ imm) ? 1 : 0 \\
16 & 10000 & R3-Type & \texttt{sll} rd = rs1 $\ll$ rs2 \\
17 & 10001 & R2I-Type & \texttt{slli} rd = rs1 $\ll$ imm \\
18 & 10010 & R3-Type & \texttt{srl} rd = rs1 $\ggg$ rs2 (logical) \\
19 & 10011 & R2I-Type & \texttt{srli} rd = rs1 $\ggg$ imm (logical) \\
20 & 10100 & R3-Type & \texttt{sra} rd = rs1 $\gg$ rs2 (arithmetic) \\
21 & 10101 & R2I-Type & \texttt{srai} rd = rs1 $\gg$ imm (arithmetic) \\
22 & 10110 & R2I-Type & \texttt{load} rd = mem[rs1 + imm] \\
23 & 10111 & R2I-Type & \texttt{store} mem[rd + imm] = rs1 \\
24 & 11000 & RI-Type & \texttt{jmp} PC = PC + rd + imm \\
25 & 11001 & R2I-Type & \texttt{beq} if rs1==rd, PC = PC + imm \\
26 & 11010 & R2I-Type & \texttt{bne} if rs1$\neq$rd, PC = PC + imm \\
27 & 11011 & R2I-Type & \texttt{blt} if rs1$<$rd, PC = PC + imm \\
28 & 11100 & R2I-Type & \texttt{bgt} if rs1$>$rd, PC = PC + imm \\
29 & 11101 & RI-Type & \texttt{end} terminate execution \\
\bottomrule
\end{longtable}

\textbf{Pattern}: Even opcodes (0, 2, 4, ..., 20) are R3-Type (register-register). Odd opcodes (1, 3, 5, ..., 21) are R2I-Type (register-immediate). Opcodes 22--28 are memory and branch instructions (all R2I-Type). Opcodes 24 and 29 are RI-Type.

% ===========================================================================
\section{Detailed Instruction Semantics}
% ===========================================================================

\subsection{Arithmetic Instructions}

\subsubsection{add / addi -- Addition}

\begin{tcolorbox}[colback=blue!5, colframe=blue!50!black]
\textbf{add} \texttt{\%rs1, \%rs2, \%rd} $\quad\Rightarrow\quad$ \texttt{rd = rs1 + rs2}

\textbf{addi} \texttt{\%rs1, imm, \%rd} $\quad\Rightarrow\quad$ \texttt{rd = rs1 + imm}
\end{tcolorbox}

Addition of two 32-bit signed integers. If the result exceeds 32 bits (overflow), the upper 32 bits are stored in \texttt{x31}. In Java, integer overflow wraps around naturally (two's complement), so the lower 32 bits are in the result and the overflow can be detected by computing the full 64-bit result using \texttt{long}:

\begin{lstlisting}[style=javastyle]
long result64 = (long)op1 + (long)op2;
int result32 = (int)result64;            // lower 32 bits
int overflow = (int)(result64 >> 32);    // upper 32 bits -> x31
\end{lstlisting}

\subsubsection{sub / subi -- Subtraction}

\begin{tcolorbox}[colback=blue!5, colframe=blue!50!black]
\textbf{sub} \texttt{\%rs1, \%rs2, \%rd} $\quad\Rightarrow\quad$ \texttt{rd = rs1 - rs2}

\textbf{subi} \texttt{\%rs1, imm, \%rd} $\quad\Rightarrow\quad$ \texttt{rd = rs1 - imm}
\end{tcolorbox}

Similar to addition but performs subtraction. Overflow handling is analogous.

\textbf{Common pattern}: \texttt{sub \%x3, \%x3, \%x3} sets \texttt{x3 = 0}. This is equivalent to \texttt{add \%x0, \%x0, \%x3} and is a standard idiom for clearing a register.

\subsubsection{mul / muli -- Multiplication}

\begin{tcolorbox}[colback=blue!5, colframe=blue!50!black]
\textbf{mul} \texttt{\%rs1, \%rs2, \%rd} $\quad\Rightarrow\quad$ \texttt{rd = (rs1 * rs2)[31:0]}, \texttt{x31 = (rs1 * rs2)[63:32]}

\textbf{muli} \texttt{\%rs1, imm, \%rd} $\quad\Rightarrow\quad$ \texttt{rd = (rs1 * imm)[31:0]}, \texttt{x31 = (rs1 * imm)[63:32]}
\end{tcolorbox}

Multiplying two 32-bit integers produces a 64-bit result. The lower 32 bits go to \texttt{rd}, and the upper 32 bits go to \texttt{x31}.

\begin{lstlisting}[style=javastyle]
long result64 = (long)op1 * (long)op2;
int rd_value = (int)result64;
int x31_value = (int)(result64 >> 32);
\end{lstlisting}

\subsubsection{div / divi -- Division}

\begin{tcolorbox}[colback=blue!5, colframe=blue!50!black]
\textbf{div} \texttt{\%rs1, \%rs2, \%rd} $\quad\Rightarrow\quad$ \texttt{rd = rs1 / rs2} (quotient), \texttt{x31 = rs1 \% rs2} (remainder)

\textbf{divi} \texttt{\%rs1, imm, \%rd} $\quad\Rightarrow\quad$ \texttt{rd = rs1 / imm}, \texttt{x31 = rs1 \% imm}
\end{tcolorbox}

Integer division. The quotient goes to \texttt{rd} and the remainder goes to \texttt{x31}. This is heavily used in the palindrome and prime test programs.

\begin{lstlisting}[style=javastyle]
int rd_value = op1 / op2;   // integer division (truncates)
int x31_value = op1 % op2;  // remainder
\end{lstlisting}

\subsection{Logical Instructions}

\subsubsection{and / andi -- Bitwise AND}

\begin{tcolorbox}[colback=green!5, colframe=green!50!black]
\texttt{rd = rs1 \& rs2} \quad or \quad \texttt{rd = rs1 \& imm}
\end{tcolorbox}

Performs bitwise AND. Each bit of the result is 1 only if the corresponding bits of both operands are 1.

\textbf{Use case}: Masking -- extracting specific bits from a value. For example, \texttt{andi \%x3, 1, \%x4} extracts the least significant bit of \texttt{x3}.

\subsubsection{or / ori -- Bitwise OR}

\begin{tcolorbox}[colback=green!5, colframe=green!50!black]
\texttt{rd = rs1 | rs2} \quad or \quad \texttt{rd = rs1 | imm}
\end{tcolorbox}

Performs bitwise OR. Each bit of the result is 1 if either corresponding bit is 1.

\textbf{Use case}: Setting specific bits. For example, \texttt{ori \%x3, 0x80, \%x3} sets bit 7 of \texttt{x3}.

\subsubsection{xor / xori -- Bitwise XOR}

\begin{tcolorbox}[colback=green!5, colframe=green!50!black]
\texttt{rd = rs1 $\oplus$ rs2} \quad or \quad \texttt{rd = rs1 $\oplus$ imm}
\end{tcolorbox}

Performs bitwise exclusive OR. The result bit is 1 if the input bits differ.

\textbf{Use case}: Toggling bits, computing parity, and encryption.

\subsection{Comparison Instructions}

\subsubsection{slt / slti -- Set Less Than}

\begin{tcolorbox}[colback=purple!5, colframe=purple!50!black]
\texttt{rd = (rs1 < rs2) ? 1 : 0} \quad or \quad \texttt{rd = (rs1 < imm) ? 1 : 0}
\end{tcolorbox}

Compares two signed integers. If \texttt{rs1} is less than \texttt{rs2} (or \texttt{imm}), the result is 1; otherwise 0.

\textbf{Use case}: Building complex control flow. For example, checking if a value is within a range.

\subsection{Shift Instructions}

\subsubsection{sll / slli -- Shift Left Logical}

\begin{tcolorbox}[colback=orange!5, colframe=orange!50!black]
\texttt{rd = rs1 << rs2} (or \texttt{<< imm})

\texttt{x31 = bits shifted out}
\end{tcolorbox}

Shifts all bits of \texttt{rs1} to the left by \texttt{rs2} (or \texttt{imm}) positions. Zeros are filled in from the right. The bits shifted out (from the left) are stored in \texttt{x31}.

\textbf{Mathematical effect}: Left shifting by $n$ is equivalent to multiplication by $2^n$.

\begin{lstlisting}[style=javastyle]
int rd_value = op1 << shiftAmount;
int x31_value = op1 >>> (32 - shiftAmount); // bits shifted out
\end{lstlisting}

\subsubsection{srl / srli -- Shift Right Logical}

\begin{tcolorbox}[colback=orange!5, colframe=orange!50!black]
\texttt{rd = rs1 >>> rs2} (logical right shift)

\texttt{x31 = bits shifted out}
\end{tcolorbox}

Shifts all bits to the right. Zeros are filled in from the left (regardless of sign). The bits shifted out (from the right) are stored in \texttt{x31}.

In Java, \texttt{>>>} is the unsigned/logical right shift operator.

\begin{lstlisting}[style=javastyle]
int rd_value = op1 >>> shiftAmount;
int x31_value = op1 << (32 - shiftAmount); // bits shifted out
\end{lstlisting}

\subsubsection{sra / srai -- Shift Right Arithmetic}

\begin{tcolorbox}[colback=orange!5, colframe=orange!50!black]
\texttt{rd = rs1 >> rs2} (arithmetic right shift, preserves sign)

\texttt{x31 = bits shifted out}
\end{tcolorbox}

Shifts all bits to the right, but fills the left with copies of the \textbf{sign bit} (MSB). This preserves the sign of negative numbers.

\textbf{Mathematical effect}: Arithmetic right shifting by $n$ is equivalent to division by $2^n$ (with rounding toward negative infinity).

In Java, \texttt{>>} is the signed/arithmetic right shift operator.

\begin{lstlisting}[style=javastyle]
int rd_value = op1 >> shiftAmount;
int x31_value = op1 << (32 - shiftAmount); // bits shifted out
\end{lstlisting}

\subsection{Memory Instructions}

\subsubsection{load -- Load Word from Memory}

\begin{tcolorbox}[colback=yellow!5, colframe=yellow!50!black]
\textbf{Format}: R2I-Type

\textbf{Assembly}: \texttt{load \%rs1, imm, \%rd}

\textbf{Semantics}: \texttt{rd = memory[rs1 + imm]}

\textbf{Opcode}: 10110 (22)
\end{tcolorbox}

Loads a 32-bit word from main memory at address \texttt{rs1 + imm} into register \texttt{rd}.

\textbf{EX stage}: Compute effective address = value of rs1 + immediate.

\textbf{MA stage}: Read \texttt{memory[effective\_address]} and pass the loaded value forward.

\textbf{RW stage}: Write the loaded value to register \texttt{rd}.

\textbf{Example}: \texttt{load \%x0, \$a, \%x4}
\begin{itemize}
    \item \texttt{rs1 = x0} (value 0)
    \item \texttt{imm = address of label a} (e.g., 0)
    \item Effective address = 0 + 0 = 0
    \item \texttt{x4 = memory[0]}
\end{itemize}

\subsubsection{store -- Store Word to Memory}

\begin{tcolorbox}[colback=yellow!5, colframe=yellow!50!black]
\textbf{Format}: R2I-Type

\textbf{Assembly}: \texttt{store \%rs1, imm, \%rd}

\textbf{Semantics}: \texttt{memory[rd + imm] = rs1}

\textbf{Opcode}: 10111 (23)
\end{tcolorbox}

Stores the value of register \texttt{rs1} into main memory at address \texttt{rd + imm}.

\textbf{Important}: Note the asymmetry with \texttt{load}. In \texttt{store}, the destination register \texttt{rd} is used as part of the \textbf{address} calculation (not as the register to write to). The value to store comes from \texttt{rs1}.

\textbf{EX stage}: Compute effective address = value of rd + immediate.

\textbf{MA stage}: Write value of rs1 to \texttt{memory[effective\_address]}.

\textbf{RW stage}: No register write needed.

\subsection{Control Flow Instructions -- Detailed}

\subsubsection{jmp -- Unconditional Jump}

\begin{tcolorbox}[colback=red!5, colframe=red!50!black]
\textbf{Format}: RI-Type

\textbf{Assembly}: \texttt{jmp label} or \texttt{jmp \%rd} (one of rd or imm is used, the other is 0)

\textbf{Semantics}: \texttt{PC = PC + rd + imm}

\textbf{Opcode}: 11000 (24)
\end{tcolorbox}

\textbf{Critical note from ISA spec}: In assembly, either \texttt{rd} or \texttt{imm} is used. In machine code, the unused one is set to zero. In the ParsedProgram, the used one is placed in the \texttt{destinationOperand} field.

The new PC is computed as: $\text{PC}_{\text{new}} = \text{PC}_{\text{current}} + \text{rd} + \text{imm}$

Since PC was already incremented during IF, the actual value of ``PC'' here refers to the PC \emph{before} increment, i.e., the address of the \texttt{jmp} instruction itself. This must be handled carefully during implementation.

\subsubsection{beq -- Branch if Equal}

\begin{tcolorbox}[colback=red!5, colframe=red!50!black]
\textbf{Format}: R2I-Type

\textbf{Assembly}: \texttt{beq \%rs1, \%rd, label}

\textbf{Semantics}: If \texttt{rs1 == rd}, then \texttt{PC = PC + imm}

\textbf{Opcode}: 11001 (25)
\end{tcolorbox}

\textbf{ParsedProgram representation}: The two registers being compared (\texttt{rs1} and \texttt{rd} in assembly notation) are placed in \texttt{sourceOperand1} and \texttt{sourceOperand2}. The label/offset is placed in \texttt{destinationOperand}.

\subsubsection{bne -- Branch if Not Equal}

\begin{tcolorbox}[colback=red!5, colframe=red!50!black]
\textbf{Semantics}: If \texttt{rs1 != rd}, then \texttt{PC = PC + imm}

\textbf{Opcode}: 11010 (26)
\end{tcolorbox}

\subsubsection{blt -- Branch if Less Than}

\begin{tcolorbox}[colback=red!5, colframe=red!50!black]
\textbf{Semantics}: If \texttt{rs1 < rd}, then \texttt{PC = PC + imm}

\textbf{Opcode}: 11011 (27)
\end{tcolorbox}

\subsubsection{bgt -- Branch if Greater Than}

\begin{tcolorbox}[colback=red!5, colframe=red!50!black]
\textbf{Semantics}: If \texttt{rs1 > rd}, then \texttt{PC = PC + imm}

\textbf{Opcode}: 11100 (28)
\end{tcolorbox}

\subsection{Special Instruction: end}

\begin{tcolorbox}[colback=gray!10, colframe=gray!50!black]
\textbf{Format}: RI-Type

\textbf{Assembly}: \texttt{end}

\textbf{Semantics}: Terminate program execution.

\textbf{Opcode}: 11101 (29)

Both \texttt{rd} and \texttt{imm} fields are unused (set to 0).
\end{tcolorbox}

When the \texttt{end} instruction reaches the RW stage, the simulator calls \texttt{Simulator.setSimulationComplete(true)} to break the simulation loop.

% ===========================================================================
\section{Object File Format}
% ===========================================================================

The ToyRISC assembler converts assembly programs into object files (binary format). The object file is what the simulator reads and loads into memory.

\subsection{Structure}

The object file contains a sequence of integers (one per line in text representation, or as binary data):

\begin{enumerate}
    \item The first integer is the \textbf{address of the first instruction} (i.e., the value to set PC to). This equals $N_d$, the number of data words.
    \item The subsequent integers are the data and code words, in order. Data words come first (addresses 0 to $N_d - 1$), followed by code words (addresses $N_d$ to $N_t - 1$).
\end{enumerate}

\subsection{Loading the Object File}

The \texttt{loadProgram()} function in \texttt{Simulator.java} must:

\begin{enumerate}
    \item Open the object file.
    \item Read the first integer: this is $N_d$, the start address of the text segment.
    \item Read all subsequent integers and store them sequentially in memory starting at address 0.
    \item Set \texttt{PC = $N_d$} (address of the first instruction).
    \item Set \texttt{x0 = 0}, \texttt{x1 = 65535}, \texttt{x2 = 65535}.
\end{enumerate}

\begin{lstlisting}[style=javastyle, caption={Loading the program (conceptual)}]
static void loadProgram(String objectFile) {
    // Read the object file
    Scanner scanner = new Scanner(new File(objectFile));

    // First value = PC (address of first instruction)
    int firstInstructionAddress = scanner.nextInt();

    // Load data and instructions into memory
    int address = 0;
    while (scanner.hasNextInt()) {
        processor.getMainMemory()
            .setWord(address, scanner.nextInt());
        address++;
    }

    // Set PC to the first instruction
    processor.getRegisterFile()
        .setProgramCounter(firstInstructionAddress);

    // Initialize special registers
    processor.getRegisterFile().setValue(0, 0);      // x0 = 0
    processor.getRegisterFile().setValue(1, 65535);   // x1 = SP
    processor.getRegisterFile().setValue(2, 65535);   // x2 = FP
}
\end{lstlisting}

% ===========================================================================
\section{Assembly Language Syntax}
% ===========================================================================

\subsection{Directives}

\begin{itemize}
    \item \texttt{.data}: Marks the beginning of the static data segment.
    \item \texttt{.text}: Marks the beginning of the code segment.
\end{itemize}

\subsection{Labels}

Labels are symbolic names for memory addresses. They make assembly programs readable.

\begin{itemize}
    \item A label is defined by placing a name followed by a colon at the beginning of a line: \texttt{loop:}
    \item Labels in the \texttt{.data} section refer to data addresses. \texttt{a:} followed by \texttt{123} means address 0 contains 123.
    \item Labels in the \texttt{.text} section refer to instruction addresses.
    \item The special label \texttt{main:} indicates the program entry point.
    \item In instructions, labels are prefixed with \texttt{\$}: \texttt{load \%x0, \$a, \%x4}
\end{itemize}

\subsection{Operand Notation}

\begin{longtable}{p{3cm}p{3cm}p{7cm}}
\toprule
\textbf{Prefix} & \textbf{Type} & \textbf{Example} \\
\midrule
\texttt{\%} & Register & \texttt{\%x3} refers to register x3 \\
\texttt{\$} & Label (address) & \texttt{\$a} refers to the address named `a' \\
(none) & Immediate & \texttt{42} is the integer constant 42 \\
\bottomrule
\end{longtable}

\subsection{Instruction Ordering}

In assembly, the destination operand is always written \textbf{last}:

\texttt{add \%x3, \%x4, \%x5} means \texttt{x5 = x3 + x4} (destination is x5, written last).

This is opposite to some other ISAs (e.g., x86, where the destination is written first).

% ===========================================================================
\section{Complete Encoding Example}
% ===========================================================================

Let us manually encode the ``Adding Two Numbers'' program:

\begin{lstlisting}[style=asmstyle, caption={Adding Two Numbers}]
.data
a:
    123
    234
.text
main:
    load %x0, $a, %x4      ; x4 = memory[x0 + 0] = 123
    addi %x0, 1, %x3       ; x3 = x0 + 1 = 1
    load %x3, $a, %x5      ; x5 = memory[x3 + 0] = memory[1] = 234
    add %x4, %x5, %x6      ; x6 = x4 + x5 = 123 + 234 = 357
    end
\end{lstlisting}

\subsection{Data Segment}

\begin{itemize}
    \item Address 0: 123
    \item Address 1: 234
\end{itemize}

$N_d = 2$ (2 data words). Instructions start at address 2.

\subsection{Instruction Encoding}

\subsubsection{Instruction 1: \texttt{load \%x0, \$a, \%x4}}

\begin{itemize}
    \item Opcode: \texttt{load} = 10110 = 22
    \item Format: R2I-Type
    \item rs1 = x0 = 00000
    \item rd = x4 = 00100
    \item imm = address of \texttt{a} = 0 = 00000000000000000 (17 bits)
\end{itemize}

Binary: \texttt{10110 00000 00100 00000000000000000}

As 32-bit integer: $10110\;00000\;00100\;0\ldots0_2 = (22 \ll 27) | (0 \ll 22) | (4 \ll 17) | 0$

$= 22 \times 2^{27} + 0 + 4 \times 2^{17} + 0 = 2952790016 + 524288 = -1341653504$ (in signed 32-bit)

\subsubsection{Instruction 2: \texttt{addi \%x0, 1, \%x3}}

\begin{itemize}
    \item Opcode: \texttt{addi} = 00001 = 1
    \item Format: R2I-Type
    \item rs1 = x0 = 00000
    \item rd = x3 = 00011
    \item imm = 1 = 00000000000000001
\end{itemize}

Binary: \texttt{00001 00000 00011 00000000000000001}

\subsubsection{Instruction 3: \texttt{load \%x3, \$a, \%x5}}

\begin{itemize}
    \item Opcode: \texttt{load} = 10110 = 22
    \item rs1 = x3 = 00011
    \item rd = x5 = 00101
    \item imm = 0
\end{itemize}

\subsubsection{Instruction 4: \texttt{add \%x4, \%x5, \%x6}}

\begin{itemize}
    \item Opcode: \texttt{add} = 00000 = 0
    \item Format: R3-Type
    \item rs1 = x4 = 00100
    \item rs2 = x5 = 00101
    \item rd = x6 = 00110
    \item unused = 000000000000
\end{itemize}

Binary: \texttt{00000 00100 00101 00110 000000000000}

\subsubsection{Instruction 5: \texttt{end}}

\begin{itemize}
    \item Opcode: \texttt{end} = 11101 = 29
    \item Format: RI-Type
    \item rd = 0, imm = 0
\end{itemize}

Binary: \texttt{11101 00000 0000000000000000000000}

\subsection{Object File}

The complete object file would contain:
\begin{verbatim}
2       (first instruction address = Nd)
123     (data: memory[0])
234     (data: memory[1])
-1341653504   (instruction: load at memory[2])
...     (instruction: addi at memory[3])
...     (instruction: load at memory[4])
...     (instruction: add at memory[5])
...     (instruction: end at memory[6])
\end{verbatim}

% ===========================================================================
\section{Instruction Decode Logic -- Implementation Theory}
% ===========================================================================

The Operand Fetch (OF) stage must decode the 32-bit instruction and extract all necessary fields. Here is a systematic approach:

\subsection{Step 1: Extract the Opcode}

\begin{lstlisting}[style=javastyle]
int instruction = IF_OF_Latch.getInstruction();
int opcode = (instruction >>> 27) & 0x1F;
\end{lstlisting}

\subsection{Step 2: Determine the Format}

\begin{lstlisting}[style=javastyle]
// R3-Type: opcodes 0,2,4,6,8,10,12,14,16,18,20
// R2I-Type: opcodes 1,3,5,7,9,11,13,15,17,19,21,22,23,25,26,27,28
// RI-Type: opcodes 24, 29
boolean isR3Type = (opcode <= 21 && opcode % 2 == 0);
boolean isRIType = (opcode == 24 || opcode == 29);
boolean isR2IType = !isR3Type && !isRIType;
\end{lstlisting}

\subsection{Step 3: Extract Fields Based on Format}

\begin{lstlisting}[style=javastyle]
if (isR3Type) {
    int rs1 = (instruction >>> 22) & 0x1F;
    int rs2 = (instruction >>> 17) & 0x1F;
    int rd  = (instruction >>> 12) & 0x1F;
    // Read register values
    int val_rs1 = registerFile.getValue(rs1);
    int val_rs2 = registerFile.getValue(rs2);
} else if (isR2IType) {
    int rs1 = (instruction >>> 22) & 0x1F;
    int rd  = (instruction >>> 17) & 0x1F;
    int imm = instruction & 0x1FFFF; // 17-bit immediate
    // Sign extend
    if ((imm & (1 << 16)) != 0)
        imm = imm | 0xFFFE0000;
    int val_rs1 = registerFile.getValue(rs1);
} else { // RI-Type
    int rd  = (instruction >>> 22) & 0x1F;
    int imm = instruction & 0x3FFFFF; // 22-bit immediate
    // Sign extend
    if ((imm & (1 << 21)) != 0)
        imm = imm | 0xFFC00000;
}
\end{lstlisting}

\subsection{Step 4: Handle Special Cases for Branches}

For conditional branches (\texttt{beq}, \texttt{bne}, \texttt{blt}, \texttt{bgt}), the ISA specification says:

\begin{quote}
``In ParsedProgram, for conditional branches, the two registers that are compared are placed in sourceOperand1 and sourceOperand2. The imm value is placed in destinationOperand.''
\end{quote}

This means:
\begin{itemize}
    \item \texttt{sourceOperand1} = value of rs1
    \item \texttt{sourceOperand2} = value of rd (not the immediate!)
    \item \texttt{destinationOperand} = the immediate offset
\end{itemize}

The OF/EX latch must be designed to carry all this information.

\vspace{2cm}
\begin{center}
\textit{End of Part 2. Continue to Part 3 for complete simulator architecture, code walkthrough, test program analysis, and implementation guide.}
\end{center}

\end{document}
